\subsection{От локальных координат к многообразиям}

В этом разделе вводится минимальный геометрический аппарат, необходимый для постановки задачи на многообразии матриц фиксированного ранга. Материал основан на лекционных и семинарских материалах по дифференциальной геометрии, в частности на курсе А. О. Иванова с мехмата МГУ, а также на курсах Фоменко и Пенского \cite{ivanov-dg,fomenko-dg,penskoy-seminars}. Мы используем только те определения и факты, которые дальше нужны для риманова градиентного спуска.

\begin{definition}
Пусть $M$ --- топологическое пространство. \textit{Картой} размерности $d$ называется пара $(U,\varphi)$, где $U\subset M$ --- открытое множество, а
\[
\varphi:U\to \varphi(U)\subset \R^d
\]
--- гомеоморфизм на открытое множество $\varphi(U)$.
\end{definition}

\begin{definition}
Набор карт $\mathcal A=\{(U_\alpha,\varphi_\alpha)\}$ называется \textit{атласом}, если множества $U_\alpha$ покрывают $M$:
\[
M=\bigcup_\alpha U_\alpha.
\]
Если все карты имеют значения в $\R^d$, то $d$ называется размерностью многообразия.
\end{definition}

\begin{definition}
Топологическое пространство $M$ называется \textit{$d$-мерным топологическим многообразием}, если у каждой точки $p\in M$ есть карта $(U,\varphi)$ размерности $d$, то есть окрестность $U$ точки $p$ гомеоморфна открытому множеству в $\R^d$.
\end{definition}

Когда две карты пересекаются, одна и та же точка получает две системы координат. Поэтому нужно понимать, как координаты переходят друг в друга. Если $U_\alpha\cap U_\beta\ne\varnothing$, то отображение перехода имеет вид
\[
\varphi_\beta\circ\varphi_\alpha^{-1}:
\varphi_\alpha(U_\alpha\cap U_\beta)
\to
\varphi_\beta(U_\alpha\cap U_\beta).
\]

\begin{definition}
Атлас называется \textit{гладким}, если все отображения перехода между картами гладкие. Пространство $M$ вместе с таким атласом называется \textit{гладким многообразием}.
\end{definition}

\begin{definition}
Пусть $M$ и $N$ --- гладкие многообразия, а $F:M\to N$ --- отображение. Оно называется \textit{гладким}, если в любых картах $(U,\varphi)$ на $M$ и $(V,\psi)$ на $N$ координатное представление
\[
\psi\circ F\circ\varphi^{-1}
\]
является гладким отображением между открытыми подмножествами евклидовых пространств.
\end{definition}

Если $F$ взаимно однозначно, гладко и обратное отображение $F^{-1}$ тоже гладко, то $F$ называется \textit{диффеоморфизмом}. В этом случае многообразия считаются одинаковыми с точки зрения гладкой геометрии.

Композиция гладких отображений снова гладкая; это обеспечивает согласованность при переходе между картами.

\subsection{Вложенные многообразия и уравнения}

В задачах оптимизации допустимое множество часто лежит в $\R^N$ и задается ограничениями. Поэтому нам прежде всего нужны вложенные многообразия.

\begin{definition}
Подмножество $M\subset\R^N$ называется \textit{вложенным гладким многообразием} размерности $d$, если для каждой точки $x\in M$ существует окрестность $W\subset\R^N$ и гладкая замена координат $\Phi:W\to\Phi(W)\subset\R^N$ с гладким обратным отображением такая, что
\[
\Phi(W\cap M)=\Phi(W)\cap\bigl(\R^d\times\{0\}^{N-d}\bigr).
\]
\end{definition}

Перед теоремой о множествах, заданных уравнениями, введем регулярные и особые точки.

\begin{definition}
Пусть $F:\R^N\to\R^k$ --- гладкое отображение. Точка $x\in\R^N$ называется \textit{регулярной точкой} отображения $F$, если дифференциал $DF(x)$ имеет максимальный ранг:
\[
\rank DF(x)=k.
\]
Если это условие не выполнено, то $x$ называется \textit{особой точкой}.
\end{definition}

\begin{definition}
Точка $y\in\R^k$ называется \textit{регулярным значением} отображения $F$, если каждая точка $x\in F^{-1}(y)$ является регулярной. Если в прообразе $F^{-1}(y)$ есть особая точка, то $y$ называется \textit{особым значением}.
\end{definition}

\begin{theorem}[Теорема о регулярном уровне]
\label{thm:regular-level}
Пусть $F:\R^N\to\R^k$ --- гладкое отображение, а $0\in\R^k$ является регулярным значением $F$. Тогда множество
\[
M=F^{-1}(0)
\]
является гладким вложенным многообразием размерности $N-k$.
\end{theorem}

\begin{proof}
Это прямое следствие теоремы о неявной функции. Условие регулярности означает, что $k$ уравнений локально независимы, поэтому $k$ координат можно выразить как гладкие функции от оставшихся $N-k$ координат. Эти свободные координаты и дают локальную карту на $M$.
\end{proof}

Обратно, вложенное многообразие локально можно задать независимой системой уравнений.

\begin{theorem}[Локальное задание вложенного многообразия уравнениями]
\label{thm:local-equations}
Пусть $M\subset\R^N$ --- вложенное гладкое многообразие размерности $d$. Тогда для каждой точки $x\in M$ существует окрестность $W\subset\R^N$ и гладкое отображение
\[
F:W\to\R^{N-d}
\]
такое, что
\[
M\cap W=F^{-1}(0),
\qquad
\rank DF(z)=N-d
\quad\text{для всех } z\in M\cap W.
\]
\end{theorem}

Доказательство этого факта мы не приводим: для нас важен вывод, что вложенные многообразия можно локально понимать как множества решений независимых гладких уравнений.

В компактной реализации будут встречаться матрицы с ортонормированными столбцами:
\[
\St(r,m)=\{U\in\R^{m\times r}:U^\top U=I_r\}.
\]
Это многообразие Stiefel. Многообразие Grassmann $\Gr(r,m)$ описывает сами $r$-мерные подпространства; в нашей задаче оно появляется через столбцовое и строковое подпространства матрицы фиксированного ранга.

\subsection{Касательные векторы и касательные пространства}

В выбранной карте $(U,\varphi)$, содержащей точку $p\in M$, касательный вектор можно сначала задать как обычный вектор
\[
v_\varphi\in\R^d.
\]
Если $(V,\psi)$ --- другая карта около $p$, координаты того же вектора пересчитываются по правилу
\begin{equation}
\label{eq:tangent-coordinate-change}
v_\psi
=
D(\psi\circ\varphi^{-1})(\varphi(p))\,v_\varphi.
\end{equation}

\begin{definition}
Касательный вектор к гладкому многообразию $M$ в точке $p$ --- это согласованный набор координатных представлений $\{v_\varphi\}$ по всем картам, содержащим $p$, причем при переходе от одной карты к другой координаты связаны формулой \eqref{eq:tangent-coordinate-change}.
\end{definition}

Множество всех касательных векторов в точке $p$ образует линейное пространство $T_pM$.

Эквивалентно, касательный вектор можно понимать как скорость гладкой кривой. Если $\gamma:(-\varepsilon,\varepsilon)\to M$ и $\gamma(0)=p$, то
\[
\left.\frac{d}{dt}\right|_{t=0}\varphi(\gamma(t))
\]
задает его координаты в карте $\varphi$.

Если $M\subset\R^N$ --- вложенное многообразие, то это описание принимает привычный вид:
\[
T_xM=
\{\dot\gamma(0):\gamma(t)\subset M,\ \gamma(0)=x\}.
\]

\begin{proposition}
\label{prop:tangent-regular-level}
Если $M=F^{-1}(0)$ --- регулярный уровень из теоремы \ref{thm:regular-level}, то
\[
T_xM=\ker DF(x).
\]
\end{proposition}

\begin{proof}
Пусть $\gamma(t)\subset M$ и $\gamma(0)=x$. Тогда $F(\gamma(t))=0$ для всех малых $t$. Дифференцируя по $t$ при $t=0$, получаем
\[
DF(x)[\dot\gamma(0)]=0.
\]
Значит, $T_xM\subset\ker DF(x)$. Обратное включение следует из локальной параметризации, полученной из теоремы о неявной функции: каждый вектор из ядра дифференциала является производной некоторой кривой на уровне $F=0$.
\end{proof}

Дальше эта схема используется для множества матриц фиксированного ранга: сначала описывается допустимое множество, затем его касательное пространство, а затем градиентный шаг.

\subsection{Риманова метрика, дифференциал и градиент}

\begin{definition}
\textit{Риманова метрика} на гладком многообразии $M$ --- это выбор скалярного произведения
\[
g_p:T_pM\times T_pM\to\R
\]
в каждой точке $p\in M$, причем для любых $\xi,\eta,\zeta\in T_pM$ и $a,b\in\R$ выполнены:
\[
g_p(a\xi+b\eta,\zeta)=ag_p(\xi,\zeta)+bg_p(\eta,\zeta),
\]
\[
g_p(\xi,\eta)=g_p(\eta,\xi),
\]
\[
g_p(\xi,\xi)>0\quad\text{при}\quad \xi\ne 0.
\]
Кроме того, метрика должна гладко зависеть от точки $p$.
\end{definition}

В этой работе используется метрика, индуцированная евклидовым скалярным произведением на пространстве матриц:
\[
g_X(\xi,\eta)=\langle \xi,\eta\rangle
=\tr(\xi^\top\eta).
\]

\begin{definition}
Пусть $f:M\to\R$ --- гладкая функция. Дифференциал функции $f$ в точке $p$ --- это линейное отображение
\[
df_p:T_pM\to\R,
\]
которое на касательном векторе $\xi=\dot\gamma(0)$ задается как
\[
df_p[\xi]=\left.\frac{d}{dt}\right|_{t=0}f(\gamma(t)).
\]
\end{definition}

\begin{definition}
Риманов градиент функции $f$ в точке $p$ --- это единственный вектор $\grad f(p)\in T_pM$, такой что
\[
g_p(\grad f(p),\xi)=df_p[\xi]
\qquad
\text{для всех } \xi\in T_pM.
\]
\end{definition}

\begin{proposition}
\label{prop:projected-gradient}
Пусть $M\subset\R^N$ --- вложенное многообразие с индуцированной евклидовой метрикой, а $\bar f$ --- гладкое продолжение функции $f$ на окрестность $M$ в $\R^N$. Тогда
\begin{equation}
\label{eq:riemannian-gradient-projection}
\grad f(x)=\Proj_{T_xM}\bigl(\nabla \bar f(x)\bigr).
\end{equation}
\end{proposition}

Нормальная компонента евклидова градиента не двигает точку вдоль многообразия, поэтому в римановом методе остается только касательная часть.

\subsection{Геодезические и ретракции}

В евклидовом пространстве градиентный шаг имеет вид
\[
x_{k+1}=x_k-\eta_k\nabla f(x_k).
\]
На многообразии касательный шаг нужно вернуть обратно на допустимое множество. Через геодезическую это записывается как
\[
x_{k+1}=\operatorname{Exp}_{x_k}\bigl(-\eta_k\grad f(x_k)\bigr).
\]
На практике вместо экспоненты используют ретракции.

\begin{definition}
\label{def:retraction}
Ретракцией называется отображение
\[
\Retr_p:T_pM\to M,
\]
которое удовлетворяет условиям
\[
\Retr_p(0_p)=p,
\qquad
D\Retr_p(0_p)=\operatorname{id}_{T_pM}.
\]
\end{definition}

Ретракция совпадает с геодезическим шагом в первом порядке малости, но обычно считается проще.

Общая схема риманова градиентного спуска имеет вид
\begin{equation}
\label{eq:rgd-general}
x_{k+1}=\Retr_{x_k}\bigl(-\eta_k\grad f(x_k)\bigr).
\end{equation}
Длина шага $\eta_k$ может быть фиксированной или подбираться процедурой backtracking line search.

\subsection{Переход к задаче заполнения матриц}

Пусть дана матрица наблюдений $Y\in\R^{m\times n}$, но известны только элементы с индексами из множества
\[
\Omega\subseteq \{1,\ldots,m\}\times\{1,\ldots,n\}.
\]
Введем маску наблюдений $W\in\R^{m\times n}$:
\[
W_{ij}=
\begin{cases}
1, & (i,j)\in\Omega,\\
0, & (i,j)\notin\Omega.
\end{cases}
\]
Тогда оператор наблюдаемых элементов записывается как
\[
\Proj_\Omega(X)=W\odot X,
\]
где $\odot$ --- адамарово, то есть покомпонентное, произведение.

Классическая выпуклая постановка:
\begin{equation}
\label{eq:convex}
\min_{X \in \R^{m\times n}}
\frac12 \fro{\Proj_\Omega(X-Y)}^2 + \lambda \nuc{X}.
\end{equation}

Геометрическая fixed-rank постановка:
\begin{equation}
\label{eq:fixed-rank}
\min_{X \in \R^{m\times n}} \frac12 \fro{\Proj_\Omega(X-Y)}^2
\qquad \text{при условии} \qquad \rank(X)=r.
\end{equation}
Оптимизация ведется по множеству
\[
\Mcal_r=\{X\in\R^{m\times n}:\rank(X)=r\}.
\]

Для шумных данных используем регуляризованную версию:
\begin{equation}
\label{eq:fixed-rank-regularized}
\min_{X \in \Mcal_r} f_\lambda(X)
= \frac12 \fro{\Proj_\Omega(X-Y)}^2 + \frac{\lambda}{2}\fro{X}^2,
\qquad \lambda\ge 0.
\end{equation}
Штраф $\fro{X}^2$ стабилизирует метод при шуме и завышенном выбранном ранге.

\subsection{Многообразие матриц фиксированного ранга}

\begin{proposition}
\label{prop:fixed-rank-manifold}
Множество
\[
\Mcal_r=\{X\in\R^{m\times n}:\rank(X)=r\}
\]
является гладким вложенным многообразием размерности
\[
\dim \Mcal_r=r(m+n-r).
\]
\end{proposition}

\begin{proof}
Рассмотрим матрицу $X\in\Mcal_r$. У нее есть невырожденный минор размера $r\times r$. Перестановкой строк и столбцов можно считать, что этот минор находится в левом верхнем углу. Тогда в некоторой окрестности $X$ любую матрицу ранга $r$ можно записать блочно:
\[
\begin{pmatrix}
A & B\\
C & D
\end{pmatrix},
\]
где $A\in\R^{r\times r}$ невырождена. Условие ранга $r$ эквивалентно равенству
\[
D=CA^{-1}B.
\]
Значит, локально свободными параметрами являются блоки $A$, $B$ и $C$. Их общее число равно
\[
r^2+r(n-r)+(m-r)r=r(m+n-r).
\]
Так как зависимый блок $D$ выражается через свободные параметры гладко, получаем локальные координаты на $\Mcal_r$.
\end{proof}

Важно различать
\[
\{X:\rank(X)=r\}
\qquad\text{и}\qquad
\{X:\rank(X)\le r\}.
\]
Первое является гладким многообразием. Второе содержит матрицы меньших рангов и в целом имеет особенности, поэтому риманов метод формулируется на страте фиксированного ранга.

Если $X\in\Mcal_r$, то удобно использовать компактное сингулярное разложение
\[
X=U\Sigma V^\top,
\]
где $U\in\R^{m\times r}$ и $V\in\R^{n\times r}$ имеют ортонормированные столбцы, а $\Sigma\in\R^{r\times r}$ невырождена.

\begin{proposition}
\label{prop:fixed-rank-tangent}
Если $X=U\Sigma V^\top\in\Mcal_r$, то касательное пространство к $\Mcal_r$ в точке $X$ имеет вид
\[
T_X\Mcal_r =
\left\{
UMV^\top + U_pV^\top + UV_p^\top :
M \in \R^{r\times r},\,
U^\top U_p = 0,\,
V^\top V_p = 0
\right\}.
\]
\end{proposition}

Три слагаемых отвечают за изменение центрального блока, столбцового подпространства и строкового подпространства.

\begin{lemma}
\label{lemma:tangent-projection}
Пусть $P_U=UU^\top$ и $P_V=VV^\top$. Тогда ортогональная проекция произвольной матрицы $Z\in\R^{m\times n}$ на $T_X\Mcal_r$ имеет вид
\begin{equation}
\label{eq:tangent-projection-fixed-rank}
\Proj_{T_X\Mcal_r}(Z)=P_UZ+ZP_V-P_UZP_V.
\end{equation}
\end{lemma}

Для ее получения разложим $Z$ по столбцовому пространству $U$ и строковому пространству $V$:
\[
Z=P_UZP_V+P_UZ(I-P_V)+(I-P_U)ZP_V+(I-P_U)Z(I-P_V).
\]
Последнее слагаемое является нормальной компонентой. Удаляя его, получаем
\[
Z-(I-P_U)Z(I-P_V)=P_UZ+ZP_V-P_UZP_V.
\]

\subsection{Градиент функции потерь на многообразии фиксированного ранга}

Рассмотрим функцию
\[
f(X)=\frac12\fro{\Proj_\Omega(X-Y)}^2.
\]

\begin{lemma}
\label{lemma:euclidean-gradient}
Евклидов градиент функции $f$ равен
\[
\nabla f(X)=\Proj_\Omega(X-Y).
\]
\end{lemma}

Действительно, для любого приращения $H\in\R^{m\times n}$ линейная часть изменения функции равна
\[
\langle \Proj_\Omega(X-Y),\Proj_\Omega(H)\rangle.
\]
Так как $\Proj_\Omega$ является ортогональной проекцией,
\[
\langle \Proj_\Omega(X-Y),\Proj_\Omega(H)\rangle
=
\langle \Proj_\Omega(X-Y),H\rangle.
\]
Отсюда следует формула для евклидова градиента.

Для регуляризованной функции \eqref{eq:fixed-rank-regularized} имеем
\[
\nabla f_\lambda(X)=\Proj_\Omega(X-Y)+\lambda X.
\]
По предложению \ref{prop:projected-gradient} риманов градиент получается проекцией этого выражения на касательное пространство:
\[
\grad_{\Mcal_r} f_\lambda(X)
=
\Proj_{T_X\Mcal_r}\bigl(\Proj_\Omega(X-Y)+\lambda X\bigr).
\]

\subsection{Ретракция для матриц фиксированного ранга}

После касательного шага нужно вернуться на $\Mcal_r$. Используем ретракцию через лучшую аппроксимацию ранга $r$:
\begin{equation}
\label{eq:fixed-rank-retraction}
R_X(\Xi)=\Pi_r(X+\Xi),
\end{equation}
где $\Pi_r$ обозначает усеченное SVD, то есть ближайшую матрицу ранга $r$ в норме Фробениуса.

\begin{proposition}
\label{prop:svd-retraction}
Отображение $R_X(\Xi)=\Pi_r(X+\Xi)$ является ретракцией на $\Mcal_r$ в окрестности $X$.
\end{proposition}

Условие $R_X(0)=X$ выполнено, так как $X$ уже имеет ранг $r$. Для касательного направления $\Xi$ матрица $X+t\Xi$ отличается от некоторой кривой на $\Mcal_r$ только на $O(t^2)$. Усеченное SVD выбирает ближайшую матрицу ранга $r$, поэтому
\[
\Pi_r(X+t\Xi)=X+t\Xi+O(t^2),
\]
то есть $DR_X(0)[\Xi]=\Xi$.

Геодезики на $\Mcal_r$ вычислять сложнее, а для градиентного метода достаточно возвращения на многообразие с правильным первым порядком. Поэтому используем ретракцию.

Итак, прямой риманов градиентный спуск имеет вид
\begin{equation}
\label{eq:rgd}
X_{k+1}
=
R_{X_k}\bigl(-\eta_k\grad_{\Mcal_r} f_\lambda(X_k)\bigr).
\end{equation}

\subsection{Компактная реализация через факторы}

Ретракция \eqref{eq:fixed-rank-retraction} требует SVD матрицы размера $m\times n$ на каждой итерации. Поэтому храним текущую точку в виде
\[
X=USV^\top,\qquad U^\top U=I_r,\quad V^\top V=I_r,
\]
где $U\in\R^{m\times r}$, $V\in\R^{n\times r}$, а $S\in\R^{r\times r}$.

Это не обычный градиентный спуск по независимым переменным $U,S,V$: шаг по-прежнему является римановым шагом по матрице $X$.

Пусть
\[
G=\nabla f_\lambda(X)=\Proj_\Omega(X-Y)+\lambda X.
\]
Тогда компоненты проекции на касательное пространство можно записать так:
\[
M=U^\top G V,
\]
\[
U_\perp=GV-UM,\qquad V_\perp=G^\top U-VM^\top.
\]
В этих обозначениях
\[
\grad_{\Mcal_r} f_\lambda(X)
=UMV^\top+U_\perp V^\top+UV_\perp^\top.
\]
Матрица $M$ меняет центральный блок, а $U_\perp$ и $V_\perp$ --- столбцовое и строковое подпространства.

Пусть длина шага равна $\eta$. Обозначим
\[
\Delta M=-\eta M,\qquad
\Delta U=-\eta U_\perp,\qquad
\Delta V=-\eta V_\perp.
\]
Тогда касательный шаг записывается как
\begin{equation}
\label{eq:compact-tangent-step}
X+\Xi
=U(S+\Delta M)V^\top+\Delta U V^\top+U\Delta V^\top.
\end{equation}

Сделаем QR-разложения
\[
\Delta U=Q_UR_U,\qquad \Delta V=Q_VR_V.
\]
Тогда
\[
X+\Xi
=
\begin{bmatrix}U & Q_U\end{bmatrix}
\begin{bmatrix}
S+\Delta M & R_V^\top\\
R_U & 0
\end{bmatrix}
\begin{bmatrix}V & Q_V\end{bmatrix}^\top.
\]
Матрица $X+\Xi$ представлена через малое ядро
\[
K=
\begin{bmatrix}
S+\Delta M & R_V^\top\\
R_U & 0
\end{bmatrix},
\]
размером не больше $2r\times 2r$. Для ретракции достаточно SVD этого ядра:
\[
K=\widetilde U\widetilde S\widetilde V^\top.
\]
После усечения до ранга $r$ получаем новые факторы:
\[
U_+=\begin{bmatrix}U & Q_U\end{bmatrix}\widetilde U_r,\qquad
S_+=\widetilde S_r,\qquad
V_+=\begin{bmatrix}V & Q_V\end{bmatrix}\widetilde V_r.
\]
Итоговая новая точка равна
\[
X_+=U_+S_+V_+^\top.
\]

\begin{algorithm}[H]
\caption{Компактный риманов градиентный спуск на $\Mcal_r$}
\label{alg:compact-rgd}
\begin{algorithmic}[1]
\Require наблюдаемая матрица $Y$, маска $\Omega$, ранг $r$, параметр $\lambda$, начальные факторы $U_0,S_0,V_0$
\Ensure приближение $\widehat X=USV^\top$ ранга $r$
\For{$k=0,1,2,\ldots$}
    \State $X_k \gets U_kS_kV_k^\top$
    \State $G_k \gets \Proj_\Omega(X_k-Y)+\lambda X_k$
    \State $M_k \gets U_k^\top G_kV_k$
    \State $U_{\perp,k} \gets G_kV_k-U_kM_k$
    \State $V_{\perp,k} \gets G_k^\top U_k-V_kM_k^\top$
    \State выбрать $\eta_k$ с помощью backtracking line search
    \State $\Delta M_k\gets-\eta_kM_k$, $\Delta U_k\gets-\eta_kU_{\perp,k}$, $\Delta V_k\gets-\eta_kV_{\perp,k}$
    \State выполнить QR-разложения $\Delta U_k=Q_UR_U$, $\Delta V_k=Q_VR_V$
    \State $K_k\gets
    \begin{bmatrix}
    S_k+\Delta M_k & R_V^\top\\
    R_U & 0
    \end{bmatrix}$
    \State вычислить SVD малого ядра $K_k=\widetilde U\widetilde S\widetilde V^\top$
    \State $U_{k+1}\gets [U_k,Q_U]\widetilde U_r$
    \State $S_{k+1}\gets \widetilde S_r$
    \State $V_{k+1}\gets [V_k,Q_V]\widetilde V_r$
    \If{выполнен критерий остановки}
        \State \textbf{break}
    \EndIf
\EndFor
\State \Return $\widehat X=U_kS_kV_k^\top$
\end{algorithmic}
\end{algorithm}

Компактная реализация сохраняет геометрию fixed-rank многообразия, но заменяет SVD большой матрицы на QR-разложения тонких матриц и SVD малого ядра.
